\documentclass[10pt]{article}
\usepackage[a4paper,landscape,margin=1.4cm,headheight=14pt]{geometry}
\usepackage{fontspec}
\usepackage[shorthands=off,brazilian]{babel}
\usepackage[dvipsnames]{xcolor}
\usepackage[normalem]{ulem}
\usepackage{paracol}
\usepackage{fancyhdr}
\usepackage{amssymb}
\definecolor{antes}{HTML}{B3261E}
\definecolor{depois}{HTML}{1E6B32}
\definecolor{cita}{HTML}{6B6B6B}
\setlength{\columnsep}{1.2cm}
\setlength{\parindent}{0pt}
\setlength{\parskip}{4pt}
\newcommand{\num}[2]{\textsuperscript{\textbf{\color{#1}[#2]}}}
\newcommand{\chgA}[2]{\num{antes}{#1}{\color{antes}#2}}
\newcommand{\chgD}[2]{\num{depois}{#1}{\color{depois}#2}}
\newcommand{\tagA}[1]{\num{antes}{#1}}
\newcommand{\tagD}[1]{\num{depois}{#1}}
\newcommand{\corte}{\textit{(corte)}}
\newcommand{\cita}[1]{{\color{cita}\small[#1]}}
\newcommand{\nota}[1]{{\color{cita}\small\textit{[#1]}}}
\newcommand{\ia}[1]{\textcolor{red}{#1}}
\newcommand{\notas}[1]{\par\vspace{2pt}{\footnotesize\color{black!75}#1}\par\vspace{4pt}\hrule\vspace{6pt}}
\pagestyle{fancy}
\fancyhf{}

\lhead{\textbf{\color{antes}Antes}}
\chead{\small 02-Desindustrialização (Teoria).tex, linhas 3-112}
\rhead{\textbf{\color{depois}Depois}}
\cfoot{\small revisao-academica-ptbr, rodada de 25/11/2026 \quad p.~\thepage}

\begin{document}
{\large\textbf{Comparação antes e depois}}\par 02-Desindustrialização (Teoria).tex, linhas 3-112\par\smallskip 49 mudanças propostas; 50 sinalizações (35 marcas de rascunho, só contadas). Legenda: \chgA{n}{texto substituído} à esquerda, \chgD{n}{texto novo} à direita, com o número de \texttt{mudancas.tsv}; \cita{chave, página} citação; \nota{nota do autor}; \sout{tachado} e \ia{vermelho}: sugestões anteriores da IA no texto.\par\medskip\hrule\medskip
\par
\begin{paracol}{2}
\switchcolumn[0]*[{\subsection*{Teoria do quê?}}]
Nas ciências sociais, \chgA{1}{raramente os problemas são prescritivos de uma
única abordagem teórica.} \chgA{2}{Distinto das ciências naturais, onde
a teoria assume quase sempre um papel de conhecimento acumulado sobre o tópico
em questão, cientistas sociais empregam a noção de teoria de forma mais ampla.} Para eles, teoria abarca tanto uma maneira particular de olhar e resolver seus problemas como uma orientação da própria identidade do pesquisador e do objetivo de sua pesquisa. \chgA{3}{É quase
inexistente para estes uma área de pesquisa onde} "teoria" signifique apenas
um conjunto relativamente bem testado e reprodutível de hipóteses sobre dado
fenômeno \nota{citar fonte}. 
\switchcolumn[1]
Nas ciências sociais, \chgD{1}{raramente um problema prescreve uma única abordagem teórica.} \chgD{2}{Diferentemente das ciências naturais, onde a teoria quase sempre é conhecimento acumulado sobre o tópico, cientistas sociais empregam a noção de forma mais ampla.} Para eles, teoria abarca tanto uma maneira particular de olhar e resolver seus problemas como uma orientação da própria identidade do pesquisador e do objetivo de sua pesquisa. \chgD{3}{Quase não há, entre eles, área de pesquisa em que} "teoria" signifique apenas
um conjunto relativamente bem testado e reprodutível de hipóteses sobre dado
fenômeno \nota{citar fonte}. 
\switchcolumn[0]*[{\notas{\textbf{[1]} \textit{R13 Ponha a ação no verbo}: "Ser prescritivo de" esconde a ação num adjetivo; com o verbo, a frase diz direto o que o problema faz.\par \textbf{[2]} \textit{R53 Termo técnico: no máximo duas vezes por parágrafo}: "teoria" aparece três vezes no parágrafo; "a noção" já a retoma sem ambiguidade. Também corrige "Distinto das" (adjetivo) para o advérbio e corta "assume um papel de" e "em questão".\par \textbf{[3]} \textit{R35 Seis cortes}: "É quase inexistente" é uma locução que cabe em "quase não há"; "em que" é a forma formal para lugar abstrato.\par \textbf{Sinalização} \textit{R53 Termo técnico: no máximo duas vezes por parágrafo} (Nas ciências sociais): Depois da mudança proposta, "teoria" fica com duas ocorrências mais a palavra entre aspas, que é citação do conceito e não conta.\par Marcas de rascunho neste parágrafo: 1 (mantidas, sem mudança).}}]
O caráter \sout{multi-paradigmático}\ia{multiparadigmático} da sociologia \nota{citar fonte} \chgA{4}{eliminou também até o momento}
a possibilidade de \sout{se quer}\ia{sequer} um consenso sobre premissas epistemológicas.
Como consequência, seções de "teoria" quase sempre representam
um conjunto de premissas, autores e "linhas teóricas" \chgA{5}{aos quais os pesquisadores
tomam predileção, com frequência
mais em função de suas trajetórias institucionais do que propriamente pela
demanda} de seus objetos de pesquisa. 
\switchcolumn[1]
O caráter \sout{multi-paradigmático}\ia{multiparadigmático} da sociologia \nota{citar fonte} \chgD{4}{eliminou até o momento}
a possibilidade de \sout{se quer}\ia{sequer} um consenso sobre premissas epistemológicas.
Como consequência, seções de "teoria" quase sempre representam
um conjunto de premissas, autores e "linhas teóricas" \chgD{5}{que os pesquisadores preferem, com frequência mais em função de suas trajetórias institucionais do que pela demanda} de seus objetos de pesquisa. 
\switchcolumn[0]*[{\notas{\textbf{[4]} \textit{R22 Use o conectivo quando o leitor não veria a relação sem ele}: O "também" é aditivo e não marca relação com o parágrafo anterior que o leitor não veja sozinho.\par \textbf{[5]} \textit{R35 Seis cortes}: "Tomar predileção a" não é regência corrente; "preferem" diz o mesmo em uma palavra. "Propriamente" não acrescenta.\par Marcas de rascunho neste parágrafo: 3 (mantidas, sem mudança).}}]
\nota{nota da IA}
\switchcolumn[1]
\nota{nota da IA}
\switchcolumn[0]*[{\notas{Marcas de rascunho neste parágrafo: 1 (mantidas, sem mudança).}}]
Essa característica \chgA{6}{implica em algumas consequências}, boas e ruins. \chgA{7}{Do ponto
de vista positivo, a discussão teórica e metodológica nas ciências sociais
costuma adquirir um caráter mais robusto.} Cientistas sociais pensam e discutem com mais
afinco sobre a natureza de seus objetos, a validade ou parcialidade de suas premissas,
o próprio lugar do pesquisador e as consequências de sua atividade \nota{citar fontes}.
\switchcolumn[1]
Essa característica \chgD{6}{implica algumas consequências}, boas e ruins. \chgD{7}{Do lado positivo, a discussão teórica e metodológica nas ciências sociais costuma ser mais robusta.} Cientistas sociais pensam e discutem com mais
afinco sobre a natureza de seus objetos, a validade ou parcialidade de suas premissas,
o próprio lugar do pesquisador e as consequências de sua atividade \nota{citar fontes}.
\switchcolumn[0]*[{\notas{\textbf{[6]} \textit{Gramática (sem regra no guia)}: Na prosa formal, "implicar" (acarretar) é transitivo direto.\par \textbf{[7]} \textit{R36 Troque a fórmula que finge profundidade pelo verbo simples}: "Adquirir um caráter mais robusto" é "ser mais robusta"; "do ponto de vista" positivo vira "do lado positivo", que é o que se contrapõe à "contrapartida" do parágrafo seguinte.\par Marcas de rascunho neste parágrafo: 1 (mantidas, sem mudança).}}]
A contrapartida é que as escolhas teóricas com frequência enrijecem a capacidade
de observação, quando não \sout{pressupões}\ia{pressupõe} predeterminam as próprias conclusões da investigação. \chgA{8}{Também
implica que aquilo que é chamado de teoria é muito distinto de um trabalho
para outro, gerando confusão ao leitor.}
Teorias de diversos níveis analíticos
se confundem e raramente o resultado de um trabalho é uma contribuição cumulativa
àquele campo de conhecimento -- ao menos não para além \chgA{9}{do grupo que
previamente já compartilhava} da mesma "linha teórica".
\switchcolumn[1]
A contrapartida é que as escolhas teóricas com frequência enrijecem a capacidade
de observação, quando não \sout{pressupões}\ia{pressupõe} predeterminam as próprias conclusões da investigação. \chgD{8}{O sentido de teoria também muda muito de um trabalho para outro, o que confunde o leitor.}
Teorias de diversos níveis analíticos
se confundem e raramente o resultado de um trabalho é uma contribuição cumulativa
àquele campo de conhecimento -- ao menos não para além \chgD{9}{do grupo que já compartilhava} da mesma "linha teórica".
\switchcolumn[0]*[{\notas{\textbf{[8]} \textit{R04 Resolva a repetição reestruturando a frase}: A frase tinha três "que" e nenhum sujeito expresso para "implica" (D9). Assumi que o sujeito é o caráter multiparadigmático; rejeite se for outro.\par \textbf{[9]} \textit{R35 Seis cortes}: "Previamente" e "já" dizem a mesma coisa.\par \textbf{Sinalização} \textit{R53 Termo técnico: no máximo duas vezes por parágrafo} (àquele campo de conhecimento): "campo" foi aprovado só no sentido de Bourdieu. Nesta seção nenhum uso é bourdieusiano (linhas 35, 75, 78, 88, 105, 108): não mexi neles, mas um leitor pode ler "campo" como conceito. Se quiser, troque os de sentido comum por "área".\par Marcas de rascunho neste parágrafo: 2 (mantidas, sem mudança).}}]
Dada a multiplicidade de perspectivas válidas para abordar um problema
e a ausência de uma condição de ciência normal, talvez não seja possível solucionar
por inteiro esta problemática. \chgA{10}{O que se pode, entretanto, é sistematizar adequadamente
algumas respostas para facilitar ao leitor o entendimento do que, para o autor,
é teoria.}
\switchcolumn[1]
Dada a multiplicidade de perspectivas válidas para abordar um problema
e a ausência de uma condição de ciência normal, talvez não seja possível solucionar
por inteiro esta problemática. \chgD{10}{Posso, entretanto, sistematizar algumas respostas para que o leitor entenda o que é teoria nesta tese.}
\switchcolumn[0]*[{\notas{\textbf{[10]} \textit{R27 Use a primeira pessoa para os atos do autor}: "O autor" em terceira pessoa destoa do singular usado no resto do texto (D3); a frase também perde um "que" e o advérbio "adequadamente", que não informa.}}]
\chgA{11}{Assim, essa seção propõe esclarecer o que significa, no escopo deste trabalho
e no contexto da discussão sobre desindustrialização, o significado de teoria. Isso
significa esclarecer:} \chgA{12}{1- o papel desta na prática sociológica adotada; 2 - quais abordagens teóricas são consideradas mais frutíferas para o entendimento do objeto; 3- o grau de conhecimento acumulado se tem sob o mesmo.}
\switchcolumn[1]
\chgD{11}{Assim, esta seção esclarece o que significa teoria no escopo deste trabalho e no contexto da discussão sobre desindustrialização. São três pontos:} \chgD{12}{(1) o papel da teoria na prática sociológica adotada; (2) as abordagens teóricas mais frutíferas para entender o objeto; (3) o grau de conhecimento acumulado sobre ele.}
\switchcolumn[0]*[{\notas{\textbf{[11]} \textit{R35 Seis cortes}: "O que significa ... o significado" e "isso significa" repetem a mesma palavra três vezes; "esta" porque a seção é a presente. Mantive o percurso, que é o da introdução do capítulo (D2).\par \textbf{[12]} \textit{R43 Coordene apenas estruturas paralelas}: Os três itens passam a ser sintagmas nominais, com numeração uniforme; o terceiro estava sem "que" ("se tem sob") e com "sob" no lugar de "sobre".}}]
O segundo e o terceiro item são discutidos nas seções abaixo. Na seção 2.2 são
destrinchadas as principais linhas teóricas -- ou \emph{abordagens teóricas} -- \ia{que} são mobilizadas como explicativas, a saber:
a sociologia econômica, as vertentes institucionalistas e as discussões sobre desenvolvimento
e dependência. \chgA{13}{Já em 2.3 é resumido o que até o momento se produziu de mais relevante}
sobre o fenômeno da desindustrialização no mundo e no Brasil. \chgA{14}{Ambas partem de alguns pressupostos sobre o objeto. Por exemplo, de que a desindustrialização é um processo estrutural que, para além de uma dimensão quantitativa, reflete uma transformação qualitativa na organização
da sociedade.} Essa transformação precisa ser observada num horizonte temporal longo e numa escala macrossociológica para ser constatada. Assim, \chgA{15}{diferencio um fechamento de fábrica, crise econômica breve ou um declínio de um setor industrial particular} do fenômeno da desindustrialização enquanto processo histórico e social.
\switchcolumn[1]
O segundo e o terceiro item são discutidos nas seções abaixo. Na seção 2.2 são
destrinchadas as principais linhas teóricas -- ou \emph{abordagens teóricas} -- \ia{que} são mobilizadas como explicativas, a saber:
a sociologia econômica, as vertentes institucionalistas e as discussões sobre desenvolvimento
e dependência. \chgD{13}{Já a seção 2.3 resume o que se produziu até o momento de mais relevante}
sobre o fenômeno da desindustrialização no mundo e no Brasil. \chgD{14}{Ambas partem de alguns pressupostos sobre o objeto, como o de que a desindustrialização é um processo estrutural: para além de uma dimensão quantitativa, ela reflete uma transformação qualitativa na organização da sociedade.} Essa transformação precisa ser observada num horizonte temporal longo e numa escala macrossociológica para ser constatada. Assim, \chgD{15}{diferencio um fechamento de fábrica, uma crise econômica breve ou o declínio de um setor industrial particular} do fenômeno da desindustrialização enquanto processo histórico e social.
\switchcolumn[0]*[{\notas{\textbf{[13]} \textit{R14 Faça do agente o sujeito}: A seção passa a ser o sujeito, como em "Na seção 2.2 são destrinchadas" logo antes, mas na voz ativa.\par \textbf{[14]} \textit{R04 Resolva a repetição reestruturando a frase}: "Por exemplo, de que..." era uma frase sem verbo principal; a fusão resolve isso e tira um "que" (D9).\par \textbf{[15]} \textit{R43 Coordene apenas estruturas paralelas}: Os três itens passam a ter artigo, como o primeiro.\par \textbf{Sinalização} \textit{R53 Termo técnico: no máximo duas vezes por parágrafo} (fenômeno da desindustrialização): Três ocorrências de "desindustrialização", todas necessárias: a última é o contraste com o fechamento de fábrica. Sem mudança.\par \textbf{Sinalização} \textit{R39 Decida se o destaque é da informação ou do autor} (Na seção 2.2 ... em 2.3): Números de seção digitados à mão. Se a ordem das seções mudar, ficam errados; \textbackslash{}ref com \textbackslash{}label resolve (a 2.3 já tem label).\par Marcas de rascunho neste parágrafo: 1 (mantidas, sem mudança).}}]
\chgA{16}{Antes, porém, cabe explicar o uso de teoria neste trabalho. Entendo teoria em dois níveis distintos, uma teoria geral e uma}
teoria substantiva \nota{fonte}. A teoria geral, \sout{discutido}\ia{discutida} na seção 2.2, se refere \sout{as}\ia{às} ao repertório teórico utilizado para análise. \chgA{17}{Corresponde as grandes divisões} internas de certa área disciplinar, que emergem, via
aglutinação ou diferenciação, pela confluência
temática mais direta -- sociologia do trabalho, sociologia do meio ambiente, \chgA{18}{sociologia da moral, etc --}
ou pela semelhança mais indireta refletida em métodos, referências bibliográficas e conceitos compartilhados -- sociologia histórica,
sociologia econômica, sociologia política, etc.
\switchcolumn[1]
\chgD{16}{Antes, porém, explico o uso de teoria neste trabalho, que entendo em dois níveis: uma teoria geral e uma}
teoria substantiva \nota{fonte}. A teoria geral, \sout{discutido}\ia{discutida} na seção 2.2, se refere \sout{as}\ia{às} ao repertório teórico utilizado para análise. \chgD{17}{Corresponde às grandes divisões} internas de certa área disciplinar, que emergem, via
aglutinação ou diferenciação, pela confluência
temática mais direta -- sociologia do trabalho, sociologia do meio ambiente, \chgD{18}{sociologia da moral, etc. --}
ou pela semelhança mais indireta refletida em métodos, referências bibliográficas e conceitos compartilhados -- sociologia histórica,
sociologia econômica, sociologia política, etc.
\switchcolumn[0]*[{\notas{\textbf{[16]} \textit{R53 Termo técnico: no máximo duas vezes por parágrafo}: "teoria" aparece cinco vezes no parágrafo; a fusão retoma a primeira por "que" e corta a moldura "cabe explicar" (D10) e "distintos", que "dois níveis" já implica.\par \textbf{[17]} \textit{R48 Crase: teste com outra preposição}: "Corresponde" pede "a"; com o artigo de "as grandes divisões", há crase (teste: "corresponde pelas").\par \textbf{[18]} \textit{Ortografia}: "Etc." leva ponto, como em "sociologia política, etc." logo abaixo (LanguageTool).\par \textbf{Sinalização} \textit{R53 Termo técnico: no máximo duas vezes por parágrafo} (sociologia do trabalho, sociologia do meio ambiente): "sociologia" aparece seis vezes, mas são nomes de subáreas numa enumeração. Repetição necessária; sem mudança.\par Marcas de rascunho neste parágrafo: 4 (mantidas, sem mudança).}}]
Essas grandes áreas \sout{sub disciplinares}\ia{subdisciplinares}
podem conter múltiplas subdivisões, \chgA{19}{as quais seria possível atribuir} certa autonomia. \chgA{20}{Para aquelas reunidas por semelhança temática, a existência e crescimento de um número de trabalhos tematicamente mais próximo do objeto justifica o enquadramento.} \chgA{21}{Conforme o número de trabalhos crescem elas podem} se subdividir, especialmente em função dos esforços e da capacidade de agregação das agendas de pesquisa dos pesquisadores.
\switchcolumn[1]
Essas grandes áreas \sout{sub disciplinares}\ia{subdisciplinares}
podem conter múltiplas subdivisões, \chgD{19}{às quais seria possível atribuir} certa autonomia. \chgD{20}{Nas áreas reunidas por semelhança temática, o enquadramento se justifica pelo número crescente de trabalhos tematicamente mais próximos do objeto.} \chgD{21}{Conforme o número de trabalhos cresce, elas podem} se subdividir, especialmente em função dos esforços e da capacidade de agregação das agendas de pesquisa dos pesquisadores.
\switchcolumn[0]*[{\notas{\textbf{[19]} \textit{R48 Crase: teste com outra preposição}: "Atribuir a" + "as quais": crase.\par \textbf{[20]} \textit{R14 Faça do agente o sujeito}: O sujeito longo ("a existência e crescimento de um número de trabalhos...") ficava longe do verbo (D10); a concordância de "próximo" passa a ser com "trabalhos".\par \textbf{[21]} \textit{R51 Concordância: procure o sujeito, e seja coerente onde há duas formas}: O núcleo do sujeito é "número"; a oração adverbial anteposta leva vírgula (R47).\par \textbf{Sinalização} \textit{R21 Desfaça as cadeias de "de"} (da capacidade de agregação das agendas de pesquisa dos pesquisadores): Cadeia de cinco "de" e "pesquisa dos pesquisadores". Uma saída: "sobretudo pelo esforço dos pesquisadores e pela capacidade de suas agendas de agregar trabalhos"; não propus porque muda um pouco o sentido.\par Marcas de rascunho neste parágrafo: 1 (mantidas, sem mudança).}}]
\chgA{22}{Já no caso daquelas de união mais indireta, onde não existe uma temática clara que enquadre os objetos, o que lhes diferencia e ajuda}
a evitar um processo de \sout{hiperespecificação}\ia{hiperespecificação} infinito
é seu grau de maturidade institucional \nota{pegar autor que embasa}. \chgA{23}{Esta é fruto também de pesquisadores que
articulam internamente a dada área maior coesão do que com aqueles "de fora".} \chgA{24}{Esta coesão não pode então ser vista como natural, sendo geralmente construída intencionalmente
enquanto um projeto/agenda de pesquisa consciente por parte de pesquisadores --} representada pelos grandes volumes "definidores" desses campos, como "Método e visão
na Sociologia Histórica" \nota{fonte} ou "Handbook of Economic Sociology" \nota{fonte} para citar apenas dois exemplos relevantes a essa pesquisa.
\switchcolumn[1]
\chgD{22}{Já nas áreas de união mais indireta, sem uma temática clara que enquadre os objetos, o que as diferencia e ajuda}
a evitar um processo de \sout{hiperespecificação}\ia{hiperespecificação} infinito
é seu grau de maturidade institucional \nota{pegar autor que embasa}. \chgD{23}{Essa maturidade resulta também da coesão maior entre os pesquisadores da área do que com os "de fora".} \chgD{24}{Essa coesão não pode, então, ser vista como natural: em geral, é construída como projeto ou agenda de pesquisa consciente dos pesquisadores --} representada pelos grandes volumes "definidores" desses campos, como "Método e visão
na Sociologia Histórica" \nota{fonte} ou "Handbook of Economic Sociology" \nota{fonte} para citar apenas dois exemplos relevantes a essa pesquisa.
\switchcolumn[0]*[{\notas{\textbf{[22]} \textit{R04 Resolva a repetição reestruturando a frase}: Um "que" a menos na frase (D9); "diferenciar" é transitivo direto ("as", não "lhes").\par \textbf{[23]} \textit{R19 Deixe claro a que cada pronome se refere}: "Esta" podia remeter a várias coisas; "essa maturidade" deixa claro. Li "articulam ... maior coesão" como construir coesão interna; rejeite se a ideia era outra (D11).\par \textbf{[24]} \textit{R35 Seis cortes}: "Intencionalmente" e "consciente" dizem o mesmo; o gerúndio "sendo construída" vira oração própria depois dos dois-pontos.\par Marcas de rascunho neste parágrafo: 4 (mantidas, sem mudança).}}]
Argumento que, em ambos os casos citados, esses grandes campos \sout{sub disciplinares}\ia{subdisciplinares} \chgA{25}{funcionam menos como uma teoria em estrito senso, acumulo de afirmações testáveis, e mais}
como "frames teóricos" \nota{fonte} que orientam o conjunto de perguntas e os mecanismos interpretativos dos fenômenos.
\chgA{26}{Sua escolha decorre tanto da elaboração do problema, como também sinaliza} um posicionamento explícito do pesquisador
enquanto parte de \sout{algum(s)}\ia{algum(ns)} campo(s) de pesquisa, além de uma tentativa de delimitar áreas de atuação privilegiada. \chgA{27}{Não são apenas utilizados, os pesquisadores aderem a estes.}
\switchcolumn[1]
Argumento que, em ambos os casos citados, esses grandes campos \sout{sub disciplinares}\ia{subdisciplinares} \chgD{25}{funcionam menos como teoria em sentido estrito, isto é, acúmulo de afirmações testáveis, e mais}
como "frames teóricos" \nota{fonte} que orientam o conjunto de perguntas e os mecanismos interpretativos dos fenômenos.
\chgD{26}{Sua escolha decorre da elaboração do problema, mas também sinaliza} um posicionamento explícito do pesquisador
enquanto parte de \sout{algum(s)}\ia{algum(ns)} campo(s) de pesquisa, além de uma tentativa de delimitar áreas de atuação privilegiada. \chgD{27}{Os pesquisadores não apenas os utilizam: aderem a eles.}
\switchcolumn[0]*[{\notas{\textbf{[25]} \textit{Ortografia}: "Estrito senso" mistura latim e português ("stricto sensu" ou "em sentido estrito"); "acúmulo" leva acento; "isto é" mostra que o aposto define a teoria em sentido estrito.\par \textbf{[26]} \textit{R43 Coordene apenas estruturas paralelas}: "Tanto ... como também" ligava um complemento ("da elaboração") a um verbo ("sinaliza"); são duas orações, ligadas por "mas também".\par \textbf{[27]} \textit{R14 Faça do agente o sujeito}: Duas orações estavam unidas só por vírgula, e o sujeito mudava no meio; com "os pesquisadores" como sujeito, a frase se lê de uma vez.\par \textbf{Sinalização} \textit{R04 Resolva a repetição reestruturando a frase} (Argumento que, em ambos os casos): Dois "que" na frase, mas os dois são necessários ("argumento que" e a relativa de "frames teóricos") e a frase se lê bem. Sem mudança.\par Marcas de rascunho neste parágrafo: 3 (mantidas, sem mudança).}}]
\chgA{28}{Como considero que o termo "abordagem" é mais intuitivo na língua portuguesa
do que frame, opto por definir estas como "abordagens teóricas".} Elas legitimam e demandam  "com o que" e "com quem" o pesquisador se propõe a trabalhar. Em outras palavras, o objeto deve ser adequável ao campo temático sob o qual aquela
abordagem reivindica legitimidade e os principais autores e conceitos precisam ser mobilizados, mesmo que criticamente ou em
discordância. Fazer uma "sociologia das emoções", por exemplo, implicaria passar obrigatoriamente por certos conceitos, autores, textos e métodos, mesmo que para rejeitá-los.
\switchcolumn[1]
\chgD{28}{Como o termo "abordagem" me parece mais intuitivo em português do que frame, chamo esses campos de "abordagens teóricas".} Elas legitimam e demandam  "com o que" e "com quem" o pesquisador se propõe a trabalhar. Em outras palavras, o objeto deve ser adequável ao campo temático sob o qual aquela
abordagem reivindica legitimidade e os principais autores e conceitos precisam ser mobilizados, mesmo que criticamente ou em
discordância. Fazer uma "sociologia das emoções", por exemplo, implicaria passar obrigatoriamente por certos conceitos, autores, textos e métodos, mesmo que para rejeitá-los.
\switchcolumn[0]*[{\notas{\textbf{[28]} \textit{R04 Resolva a repetição reestruturando a frase}: Um "que" a menos (D9) e "estas" sem referente claro trocado por "esses campos", o sujeito do parágrafo anterior.\par \textbf{Sinalização} \textit{Ortografia} (do que frame): "frame" é palavra estrangeira e pede itálico (\textbackslash{}emph\{frame\}). Não pus porque acrescentar um comando é mexer em conteúdo protegido.}}]
\chgA{29}{Dentro das abordagens teóricas com maturidade institucional na sociologia, não existe nenhuma área que possa ser entendida} como uma
"sociologia da desindustrialização". \chgA{30}{O mais próximo que posso chegar pelo critério temático seria a sociologia do desenvolvimento, mas esta} encontra-se em significativa paralisia nas últimas décadas \nota{fonte}, o que tornaria seu uso
pouco produtivo. Opto assim por uma identificação menos temática. Entre as abordagens possíveis, diversas parecem ser importantes e contemplar o fenômeno da desindustrialização, 
mas há diferenças. Utilizo então um duplo critério, sabidamente subjetivo: a abordagem teórica deve possuir um conjunto de pesquisadores
atuantes e produtivos, para propiciar um diálogo vivo e uma contribuição maior da tese, e \chgA{31}{deve encapsular sem
significativos esforços a "desindustrialização" como um objeto ao qual tem legitimidade de discurso.} Aquela que aparenta maior adequação 
e \chgA{32}{dispõe de um campo recente relativamente produtivo, é a sociologia econômica.} \chgA{33}{Por obviedade isso implica
que se entende a sociologia econômica como uma abordagem teórica parte da disciplina sociologia e não como uma teoria específica,
uma disciplina a parte ou meramente um campo temático.}
\switchcolumn[1]
\chgD{29}{Entre as abordagens teóricas com maturidade institucional na sociologia, nenhuma pode ser entendida} como uma
"sociologia da desindustrialização". \chgD{30}{Pelo critério temático, a mais próxima seria a sociologia do desenvolvimento, mas ela} encontra-se em significativa paralisia nas últimas décadas \nota{fonte}, o que tornaria seu uso
pouco produtivo. Opto assim por uma identificação menos temática. Entre as abordagens possíveis, diversas parecem ser importantes e contemplar o fenômeno da desindustrialização, 
mas há diferenças. Utilizo então um duplo critério, sabidamente subjetivo: a abordagem teórica deve possuir um conjunto de pesquisadores
atuantes e produtivos, para propiciar um diálogo vivo e uma contribuição maior da tese, e \chgD{31}{deve incorporar sem grande esforço a "desindustrialização" como um objeto sobre o qual tem legitimidade para falar.} Aquela que aparenta maior adequação 
e \chgD{32}{dispõe de um campo recente relativamente produtivo é a sociologia econômica.} \chgD{33}{Isso implica entender a sociologia econômica como uma abordagem teórica dentro da disciplina sociologia, e não como uma teoria específica, uma disciplina à parte ou um mero campo temático.}
\switchcolumn[0]*[{\notas{\textbf{[29]} \textit{R35 Seis cortes}: "Não existe nenhuma área que possa" diz em seis palavras o que "nenhuma pode" diz em duas, e tira um "que" (D9).\par \textbf{[30]} \textit{R04 Resolva a repetição reestruturando a frase}: Um "que" a menos (D9); "esta" passa a "ela", sem ambiguidade.\par \textbf{[31]} \textit{R33 Faça o teste da fala}: "Encapsular sem significativos esforços" e "legitimidade de discurso" não se diriam a um colega; o sentido fica o mesmo.\par \textbf{[32]} \textit{R45 Não separe por vírgula sujeito e verbo, nem verbo e complemento}: A vírgula separava o sujeito ("Aquela que aparenta...") do verbo "é".\par \textbf{[33]} \textit{R32 Não anuncie a importância nem promova o conceito}: "Por obviedade" só anuncia; a frase perde um "que" (D9) e corrige "à parte" (crase).\par \textbf{Sinalização} \textit{R04 Resolva a repetição reestruturando a frase} (Dentro das abordagens teóricas): Com as mudanças propostas, o parágrafo cai de cinco para dois ou três "que".\par Marcas de rascunho neste parágrafo: 1 (mantidas, sem mudança).}}]
\chgA{34}{A seção seguinte desenvolverá esta justificativa e apontará, dentro desta escolha e em abordagens relacionadas, conceitos e autores que serão mobilizados, aqui por vezes tratados como vertentes internas a abordagem teórica mobilizada.} Essas vertentes disputam espaço internamente e, no limite, poderiam adquirir autonomia suficiente para operar como abordagens teóricas próprias. \chgA{35}{Esse limite é evidentemente "embaçado"}, como qualquer categoria sociológica, e com frequência levantará disputas como a sobre se \sout{esta}\ia{está} se fazendo uma "sociologia econômica de viés \sout{bourdiesiano}\ia{bourdieusiano}", uma "sociologia \sout{bourdiesiana}\ia{bourdieusiana} sobre a economia" ou uma "sociologia do mercado de arte \sout{bourdiesiana}\ia{bourdieusiana}", para limitar a poucas as possíveis classificações. Aqui não argumento ter feito uma taxonomia adequada da diversidade de \chgA{36}{qualquer disciplina ou sub-áreas, apenas que} \emph{para fins desse trabalho} \chgA{37}{o que chamamos de teoria} é a junção de uma abordagem teórica denominada sociologia econômica, com certas vertentes desta sendo privilegiadas, e uma teoria substantiva sobre desindustrialização, que abriga qualquer produção relevante sobre o tema que possa ser mobilizada enquanto conhecimento acumulado verificável.
\switchcolumn[1]
\chgD{34}{A seção seguinte desenvolve esta justificativa e aponta, nessa escolha e em abordagens relacionadas, os conceitos e autores que mobilizo, por vezes tratados aqui como vertentes internas à abordagem.} Essas vertentes disputam espaço internamente e, no limite, poderiam adquirir autonomia suficiente para operar como abordagens teóricas próprias. \chgD{35}{Esse limite é "embaçado"}, como qualquer categoria sociológica, e com frequência levantará disputas como a sobre se \sout{esta}\ia{está} se fazendo uma "sociologia econômica de viés \sout{bourdiesiano}\ia{bourdieusiano}", uma "sociologia \sout{bourdiesiana}\ia{bourdieusiana} sobre a economia" ou uma "sociologia do mercado de arte \sout{bourdiesiana}\ia{bourdieusiana}", para limitar a poucas as possíveis classificações. Aqui não argumento ter feito uma taxonomia adequada da diversidade de \chgD{36}{qualquer disciplina ou subárea, apenas que} \emph{para fins desse trabalho} \chgD{37}{o que chamo de teoria} é a junção de uma abordagem teórica denominada sociologia econômica, com certas vertentes desta sendo privilegiadas, e uma teoria substantiva sobre desindustrialização, que abriga qualquer produção relevante sobre o tema que possa ser mobilizada enquanto conhecimento acumulado verificável.
\switchcolumn[0]*[{\notas{\textbf{[34]} \textit{R53 Termo técnico: no máximo duas vezes por parágrafo}: "abordagens teóricas" aparece três vezes no parágrafo; "à abordagem" retoma a anterior. Também corta a repetição "mobilizados ... mobilizada", passa à primeira pessoa (D3) e põe a crase em "à abordagem".\par \textbf{[35]} \textit{R30 Dispense os intensificadores}: "Evidentemente" não informa (D5); as aspas já marcam o limite como impreciso.\par \textbf{[36]} \textit{Ortografia}: "Subárea" não leva hífen; singular para concordar com "qualquer".\par \textbf{[37]} \textit{R27 Use a primeira pessoa para os atos do autor}: O resto do parágrafo está no singular ("argumento"); o plural aqui é a mistura que o perfil sinaliza (D3).\par \textbf{Sinalização} \textit{R08 Um parágrafo, uma ideia, enunciada no início} (A seção seguinte desenvolverá): Duas ideias: o anúncio da seção seguinte com a fronteira "embaçada" entre vertentes, e a definição do que é teoria nesta tese. A segunda é a conclusão da seção inteira e talvez mereça parágrafo próprio. A última frase tem quatro "que" (D9).\par \textbf{Sinalização} \textit{R04 Resolva a repetição reestruturando a frase} (como a sobre se \textbackslash{}sout\{esta\}\textbackslash{}aiflag\{está\} se fazendo): "Como a sobre se está se fazendo" é difícil de ler. Sugestão: "como a de saber se se faz uma...". Não propus porque o trecho contém marcas de rascunho.\par Marcas de rascunho neste parágrafo: 4 (mantidas, sem mudança).}}]
\chgA{38}{Vale pontuar, outras duas abordagens teóricas identificáveis, a sociologia histórica e a sociologia do trabalho, são observadas com maior frequência, 
na literatura sociológica sobre desindustrialização-- a teoria substantiva --, mas em última instância foram descartadas por apresentar menor dinamismo recente, 
no caso da primeira, e por limitar demais as características relevantes do fenômeno a serem examinadas, no caso
da segunda.} \chgA{39}{Nem por isso deixaram de ser mobilizadas em menor grau, juntamente da história econômica, de vertentes institucionalistas da ciência política e economia além da sociologia política, quando consideradas relevantes.}
\switchcolumn[1]
\chgD{38}{Outras duas abordagens teóricas, a sociologia histórica e a sociologia do trabalho, aparecem com maior frequência na literatura sociológica sobre desindustrialização -- a teoria substantiva --, mas foram descartadas: a primeira por apresentar menor dinamismo recente; a segunda por limitar demais as características relevantes do fenômeno a examinar.} \chgD{39}{Nem por isso deixaram de ser mobilizadas, em menor grau e quando relevantes, ao lado da história econômica, de vertentes institucionalistas da ciência política e da economia, e da sociologia política.}
\switchcolumn[0]*[{\notas{\textbf{[38]} \textit{R32 Não anuncie a importância nem promova o conceito}: "Vale pontuar" só anuncia (D10); sai a vírgula entre verbo e complemento (R45), e os dois motivos ficam lado a lado depois dos dois-pontos, cada um com o seu caso. "Em última instância" e "identificáveis" não acrescentam.\par \textbf{[39]} \textit{R43 Coordene apenas estruturas paralelas}: Os três itens passam a vir todos com "de" + artigo, e as duas condições ("em menor grau", "quando relevantes") ficam juntas, antes da lista.}}]
Sobre esta última, Reis \nota{colocar fonte} considera que a sociologia do desenvolvimento brasileira \chgA{40}{constituia na verdade uma} sociologia política. \chgA{41}{Independente disso, me parece inadequado enquadrar processo de desindustrialização, e do desenvolvimento em geral,} primordialmente na sociologia política. \chgA{42}{Primeiramente, porque a maioria dos referenciais que tratam da temática não são oriundos ou dialogam com esta. Em segundo lugar, pois entendo} o fenômeno da desindustrialização como sobretudo econômico. Por fim, Nahoum \nota{citar fonte} argumenta ser a sociologia econômica também devedora da sociologia do desenvolvimento, e do trabalho. Cita inclusive a contribuição de autores como Fernandes (fonte), Guerreiro Ramos (fonte) e \chgA{43}{Cardoso e Falleto} (fonte) como precursores da sociologia econômica no Brasil (fonte nahoum). Entretanto, essa dupla "herança" fortalece a perspectiva deste trabalho \chgA{44}{da necessidade uma sociologia econômico-política para tratar o fenômeno, como será defendido adiante.}
\switchcolumn[1]
Sobre esta última, Reis \nota{colocar fonte} considera que a sociologia do desenvolvimento brasileira \chgD{40}{constituía, na verdade, uma} sociologia política. \chgD{41}{Independentemente disso, parece-me inadequado enquadrar o processo de desindustrialização, e o desenvolvimento em geral,} primordialmente na sociologia política. \chgD{42}{Primeiramente, porque a maioria dos referenciais sobre o tema não se origina dela nem dialoga com ela. Em segundo lugar, porque entendo} o fenômeno da desindustrialização como sobretudo econômico. Por fim, Nahoum \nota{citar fonte} argumenta ser a sociologia econômica também devedora da sociologia do desenvolvimento, e do trabalho. Cita inclusive a contribuição de autores como Fernandes (fonte), Guerreiro Ramos (fonte) e \chgD{43}{Cardoso e Faletto} (fonte) como precursores da sociologia econômica no Brasil (fonte nahoum). Entretanto, essa dupla "herança" fortalece a perspectiva deste trabalho \chgD{44}{da necessidade de uma sociologia econômico-política para tratar o fenômeno, como defendo adiante.}
\switchcolumn[0]*[{\notas{\textbf{[40]} \textit{Ortografia}: Acento; "na verdade" no meio da oração vai entre vírgulas (R47).\par \textbf{[41]} \textit{Gramática (sem regra no guia)}: Advérbio ("independentemente"), sem pronome átono no início da frase, artigo em "o processo" e "o desenvolvimento" como segundo objeto de "enquadrar".\par \textbf{[42]} \textit{R43 Coordene apenas estruturas paralelas}: "Não são oriundos ou dialogam" negava só o primeiro termo; "nem" nega os dois. "Porque" nas duas razões mantém o paralelo.\par \textbf{[43]} \textit{Ortografia}: O sobrenome é Faletto (Enzo Faletto).\par \textbf{[44]} \textit{Gramática (sem regra no guia)}: Falta o "de" de "necessidade de"; "como defendo" segue a primeira pessoa (D3).\par \textbf{Sinalização} \textit{R53 Termo técnico: no máximo duas vezes por parágrafo} (Sobre esta última): "sociologia" sete vezes, quase todas como nome de subárea (do desenvolvimento, política, econômica). Sem mudança.\par \textbf{Sinalização} \textit{R05 Cada conectivo marca uma relação, uma vez} (Entretanto, essa dupla "herança"): "Entretanto" marca oposição, mas a frase não contraria a anterior: a dupla herança reforça o argumento. "Assim" ou "Essa dupla herança, portanto," seria a relação real. Não mudei porque depende do que você quis contrapor.\par Marcas de rascunho neste parágrafo: 7 (mantidas, sem mudança).}}]
Já \ia{o} segundo item, discutido na seção 2.3, \chgA{45}{refere-se diretamente ao conjunto de trabalhos que investigam diretamente}
o fenômeno da desindustrialização, sendo pertencentes a distintas abordagens teóricas. \chgA{46}{Aqui duas escolhas
foram feitas. Primeiramente, dividiu-se o tema entre os trabalhos nacionais, sobre especificamente o
fenômeno brasileiro, e os internacionais.} \chgA{47}{Segundo, dentro dos trabalhos nacionais a enorme
maioria encontrada se originou do campo da economia, o que dado o baixo número geral de trabalhos
sobre o tema, levou a sua incorporação.} Para a literatura internacional, o número maior
de trabalhos permitiu -- e exigiu -- um recorte mais específico. \chgA{48}{Foram priorizadas então obras oriundas
da sociologia e da história, considerados os dois campos mais próximos desta tese.} Assim, esse segundo uso de "teoria" é mais restrito e
mais próximo da descrição de \sout{teria}\ia{teoria} na metodologia científica tradicional. Ele permite elencar certo conjunto
\chgA{49}{de afirmações sobre o fenômeno ao qual seja possível atribuir certo grau de confiança e possam
em parte ser verificadas por meio da pesquisa empírica, fortalecendo ou enfraquecendo a confiança nas mesmas.}
\switchcolumn[1]
Já \ia{o} segundo item, discutido na seção 2.3, \chgD{45}{refere-se ao conjunto de trabalhos que investigam diretamente}
o fenômeno da desindustrialização, sendo pertencentes a distintas abordagens teóricas. \chgD{46}{Aqui fiz duas escolhas. Primeiramente, dividi o tema entre os trabalhos nacionais, especificamente sobre o fenômeno brasileiro, e os internacionais.} \chgD{47}{Segundo, a maioria dos trabalhos nacionais encontrados se originou do campo da economia, e, dado o baixo número geral de trabalhos sobre o tema, decidi incorporá-los.} Para a literatura internacional, o número maior
de trabalhos permitiu -- e exigiu -- um recorte mais específico. \chgD{48}{Priorizei então obras da sociologia e da história, que considero os dois campos mais próximos desta tese.} Assim, esse segundo uso de "teoria" é mais restrito e
mais próximo da descrição de \sout{teria}\ia{teoria} na metodologia científica tradicional. Ele permite elencar certo conjunto
\chgD{49}{de afirmações sobre o fenômeno às quais se pode atribuir certo grau de confiança. A pesquisa empírica pode verificá-las em parte, fortalecendo ou enfraquecendo essa confiança.}
\switchcolumn[0]*[{\notas{\textbf{[45]} \textit{R35 Seis cortes}: "Diretamente" duas vezes na mesma frase; fica o que qualifica a investigação.\par \textbf{[46]} \textit{R27 Use a primeira pessoa para os atos do autor}: As escolhas são do autor; a passiva e o "se" impessoal escondem quem age (D3, D10).\par \textbf{[47]} \textit{R30 Dispense os intensificadores}: "Enorme" sem medida sai (D5 ampliada); "o que dado ... levou a sua incorporação" vira oração com agente (D10). Rejeite se "enorme" tem um número por trás.\par \textbf{[48]} \textit{R27 Use a primeira pessoa para os atos do autor}: A escolha é do autor (D3); "considerados" concordava com "campos" de longe.\par \textbf{[49]} \textit{R51 Concordância: procure o sujeito, e seja coerente onde há duas formas}: "Ao qual" concordava com "fenômeno", mas o que recebe confiança são as "afirmações"; a frase se divide em duas e "nas mesmas" vira "essa confiança".\par \textbf{Sinalização} \textit{R09 Enuncie cada tema uma vez; depois, só o que o tópico acrescenta} (Já \textbackslash{}aiflag\{o\} segundo item, discutido na seção 2.3): Na lista das linhas 44-46, a seção 2.3 é o terceiro item (conhecimento acumulado), não o segundo. Se "segundo" aqui quer dizer o segundo nível de teoria (a teoria substantiva, linhas 55-56), vale dizer isso; se é o item da lista, é "terceiro".\par \textbf{Sinalização} \textit{R27 Use a primeira pessoa para os atos do autor} (esse segundo uso de "teoria"): Confirma a leitura acima: "segundo uso" = teoria substantiva. Sugiro "Já a teoria substantiva, discutida na seção 2.3, ..." no início do parágrafo.\par Marcas de rascunho neste parágrafo: 2 (mantidas, sem mudança).}}]
\end{paracol}
\end{document}
